\documentclass[11pt]{article}
\usepackage{fontspec}
\usepackage[french,english]{babel}
\usepackage{amsmath}
\usepackage{amssymb}
\usepackage{amsthm}
\usepackage{mathrsfs}
\usepackage{tikz-cd}
\usepackage{enumitem}
\usetikzlibrary{babel}

\newcommand{\printedpage}[1]{\marginpar{\footnotesize p.~#1}}
\newcommand{\msp}[1]{\marginpar{\footnotesize [#1]}}
\newcommand{\Hot}{\mathrm{Hot}}
\newcommand{\HOT}{\mathrm{HOT}}
\newcommand{\Cat}{\mathrm{Cat}}
\newcommand{\cX}{\mathcal{X}}

\begin{document}

\printedpage{83}
% head: A. Grothendieck : Dérivateurs, Ch. XII, Édition des Universités de Montpellier II et Paris VII
by the known computation of the functor $\pi_!$, left adjoint to the functor $\pi^*$
\[
\HOT_W(e) \overset{\text{def}}{=} \Hot_W \longrightarrow \HOT_W(\Phi)
\]
(\footnote[84]{This assumes that we have W(2), not only W(2 bis).}). Now for every object $I$ in Cat, with absolute projection $p_I : I \longrightarrow e$, the functor
\[
p_{I\,!} : \HOT_W(I) \longrightarrow \HOT_W(e) = \Hot_W
\]
must be regarded as the ``homotopic'' equivalent of the direct limit functor $\varinjlim\limits_{I}$ on diagrams of type $I^{\circ}$ in the category $\Hot_W$ - which direct limits \textit{do not exist} in $\Hot_W$ (except
\msp{page 123}
for very special categories - the only ones I know are the discrete categories (which perhaps even have to be assumed finite as well)).

Let us note that we have a diagram of canonical functors
\[
\begin{tikzcd}[column sep=tiny, row sep=small, ampersand replacement=\&]
 \& X_0 \arrow[dl, "\alpha_1"'] \arrow[dr, "\alpha_2"] \& \\
X_1 \arrow[dr, "i_1"'] \& {\shortstack{\scriptsize not\\ \scriptsize commuta-\\ \scriptsize tive!}} \& X_2 \arrow[dl, "i_2"] \\
 \& \cX, \&
\end{tikzcd}
\tag*{$(*)$}
\]
which diagram has nothing commutative about it in Cat. \textit{It is however commutative in} $\Hot_W$ (whatever the basic localizer $W$), and even in Cat\,hot, for if we denote by
\[
i_0 : X_0 \longrightarrow \cX
\]
the canonical inclusion, we have canonical homomorphisms
\[
i_1\alpha_1 \longrightarrow i_0 \longleftarrow i_2\alpha_2,
\]
which shows that $i_1\alpha_1$ and $i_2\alpha_2$ are homotopic, hence define the same arrow in $\Hot_W$.

\msp{page 124}
Let us consider in $\Phi$ the two closed subsets isomorphic to $\boldsymbol{\Delta}^1$
\[
\begin{aligned}
F_1 &= 1\backslash\Phi &&= \{0,1\}\\
F_2 &= 2\backslash\Phi &&= \{0,2\},
\end{aligned}
\]
we have
\[
F_1\cap F_2 = 0\backslash\Phi = \{0\}, \qquad F_1\cup F_2 = \Phi.
\]

\printedpage{84}
% head: A. Grothendieck : Dérivateurs, Ch. XII, Édition des Universités de Montpellier II et Paris VII
Let
\[
\begin{aligned}
1\backslash\cX &= \cX\times_{\Phi} 1\backslash\Phi &&= \cX\times_{\Phi} F_1\\
2\backslash\cX &= \cX\times_{\Phi} 2\backslash\Phi &&= \cX\times_{\Phi} F_2
\end{aligned}
\]
be the inverse images of $F_1$, $F_2$ in $\cX$, we therefore have
\[
1\backslash\cX\cap 2\backslash\cX = X_0, \qquad 1\backslash\cX\cup 2\backslash\cX = \cX,
\]
so the diagram
\[
\begin{tikzcd}[column sep=small, row sep=small, ampersand replacement=\&]
 \& X_0 \arrow[dl, hook', "\alpha'_1"'] \arrow[dr, hook, "\alpha'_2"] \& \\
1\backslash\cX \arrow[dr, hook, "i'_1"'] \& \& 2\backslash\cX \arrow[dl, hook', "i'_2"] \\
 \& \cX, \&
\end{tikzcd}
\tag*{$(**)$}
\]
deduced from the similar diagram
\[
\begin{tikzcd}[column sep=small, row sep=small, ampersand replacement=\&]
 \& \{0\} \arrow[dl, hook'] \arrow[dr, hook] \& \\
F_1 \arrow[dr, hook] \& \& F_2 \arrow[dl, hook'] \\
 \& \Phi, \&
\end{tikzcd}
\]
is (like the latter) a square of closed immersions both cartesian and cocartesian. On the other hand, we have a canonical homomorphism
\msp{page 125}
from the diagram $(*)$ to $(**)$, which is the identity on the vertices $X_0$ and $\cX$, and which on $X_1$, $X_2$ is given by the canonical inclusions
\[
\underbrace{X_1}_{\text{\scriptsize $\cX_1$ fiber of $\cX$ at 1}} \longrightarrow 1\backslash\cX, \qquad
\underbrace{X_2}_{\text{\scriptsize $\cX_2$ fiber of $\cX$ at 2}} \longrightarrow 2\backslash\cX.
\]
These inclusions, as we know, are $W$-aspheric (they even have a right adjoint), a fortiori they are in $W$. Thus the noncommutative diagram $(*)$ is isomorphic, as a diagram in $\Hot_W$, to the commutative cartesian and cocartesian diagram $(**)$ formed of closed immersions. Since $\cX$ is indeed an amalgamated sum of its two closed subsets % [sic] fermées -> fermés, cf. « les deux fermés » p. 83
$1\backslash\cX$, $2\backslash\cX$ under their intersection $X_0$, this supports the intuition which sees $\cX$ as a ``homotopic amalgamated sum''. Let us give at once the following application:

\textbf{Proposition 13.} \textit{Let there be a diagram} $\begin{tikzcd}[column sep=small, row sep=tiny, ampersand replacement=\&, baseline=(\tikzcdmatrixname-2-1.base)]
 \& X_1 \\
X_0 \arrow[ur, "\alpha_1"] \arrow[dr, "\alpha_2"'] \& \\
 \& X_2
\end{tikzcd}$ \textit{in} Cat\textit{, whence the diagram} $(*)$\textit{, p. 123,}

\printedpage{85}
% head: A. Grothendieck : Dérivateurs, Ch. XII, Édition des Universités de Montpellier II et Paris VII
\[
\begin{tikzcd}[column sep=small, row sep=small, ampersand replacement=\&]
 \& X_0 \arrow[dl, "\alpha_1"'] \arrow[dr, "\alpha_2"] \& \\
X_1 \arrow[dr, "i_1"'] \& \& X_2 \arrow[dl, "i_2"] \\
 \& \cX, \&
\end{tikzcd}
\]
\msp{page 126}
\textit{commutative in} $\Hot_W$ \textit{(but not in} Cat\textit{, unless $X_0 = \emptyset$). Let us suppose that $W$ satisfies} W(7 bis) (cocartesian square axiom, cf. p. 75). \textit{Then we have}
\[
\left\{
\begin{aligned}
\alpha_1\in W &\iff i_2\in W\\
\alpha_2\in W &\iff i_1\in W. \text{ (\footnotemark[85])}
\end{aligned}
\right.
\]
\footnotetext[85]{\textbf{NB} For the implication $\Longrightarrow$ we need only the direct form of W(7 bis), i.e. (with the notations of XVI) of Loc(6$'$a).}
\textsc{Proof}. Using the preceding homomorphism from $(*)$ to $(**)$, we see that each of the arrows $\alpha_1$, $\alpha_2$, $i_1$, $i_2$ is in $W$, if and only if the corresponding arrow $\alpha'_1$, $\alpha'_2$, $i'_1$ or $i'_2$ is. Thus the above equivalences are equivalent to
\[
\begin{aligned}
\alpha'_1\in W &\iff i'_2\in W\\
\alpha'_2\in W &\iff i'_1\in W,
\end{aligned}
\]
which are true by virtue of W(7 bis) (dual formulation) applied to $\cX$ and to the two closed subsets $1\backslash\cX$, $2\backslash\cX$ of $\cX$.

\textbf{Corollary:} \textit{Let us consider} two \textit{diagrams of type} $\Phi^{\circ}$\textit{, and a homomorphism} $f = (f_0, f_1, f_2)$ \textit{from one to the other (cf. the diagram in solid lines opposite} \texttt{[below]}\textit{):}
\[
\begin{tikzcd}[column sep=small, row sep=small, ampersand replacement=\&]
 \& X_1 \arrow[dd, "f_1"'] \arrow[drr, dotted, "i_1"] \& \& \\
X_0 \arrow[ur, "\alpha_1"] \arrow[drr, "\alpha_2"] \arrow[dd, "f_0"'] \& \& \& \cX \arrow[dd, dotted, "f_{\cX}"] \\
 \& X'_1 \arrow[drr, dotted, "i'_1"] \& X_2 \arrow[ur, dotted, "i_2"] \arrow[dd, "f_2"'] \& \\
X'_0 \arrow[ur, "\alpha'_1"] \arrow[drr, "\alpha'_2"] \& \& \& \cX', \\
 \& \& X'_2 \arrow[ur, dotted, "i'_2"] \&
\end{tikzcd}
\]
\textit{whence a homomorphism}
\msp{page 127}
\[
f_{\cX} = \int f : \quad \cX = \int
\begin{tikzcd}[column sep=small, row sep=tiny, ampersand replacement=\&, baseline=(\tikzcdmatrixname-2-1.base)]
 \& X_1 \\
X_0 \arrow[ur] \arrow[dr] \& \\
 \& X_2
\end{tikzcd}
\qquad\longrightarrow\qquad
\cX' = \int
\begin{tikzcd}[column sep=small, row sep=tiny, ampersand replacement=\&, baseline=(\tikzcdmatrixname-2-1.base)]
 \& X'_1 \\
X'_0 \arrow[ur] \arrow[dr] \& \\
 \& X'_2,
\end{tikzcd}
\]

\printedpage{86}
% head: A. Grothendieck : Dérivateurs, Ch. XII, Édition des Universités de Montpellier II et Paris VII
\textit{fitting into a cubical diagram (cf. the preceding diagram completed by dotted lines). Let us consider the following conditions.}

\begin{itemize}
\item[\textit{a)}] $f_0\in W$\textit{.}
\item[\textit{b)}] $f_1$ \textit{and} $f_2$ \textit{are} \texttt{[in]} $W$\textit{.}
\item[\textit{c)}] $f_{\cX} = \int(f_0, f_1, f_2)\in W$\textit{.}
\end{itemize}

\textit{Let us suppose the axiom} W(7) (``5 lemma'', cf. page 76) \textit{is satisfied. Then two among the conditions a), b), c) imply the third}
\[
\left\{
\begin{aligned}
a \mathbin{\&} b &\quad\Longrightarrow\quad c\\
b \mathbin{\&} c &\quad\Longrightarrow\quad a\\
c \mathbin{\&} a &\quad\Longrightarrow\quad b.
\end{aligned}
\right.
\]

Indeed, we have a homomorphism from the cubical diagram above to the similar diagram, where $X_1$, $X_2$, $X'_1$, $X'_2$ are replaced respectively not the closed subsets % [sic] pas -> par, misprint
$1\backslash\cX$, $2\backslash\cX$ of $\cX$, and $1\backslash\cX'$, $2\backslash\cX'$ of $\cX'$, the two squares, upper and lower, being inclusions of open subsets \texttt{[rather of closed subsets]} of $\cX$ resp. of $\cX'$. Thus, conditions a), b), c) are equivalent respectively to the similar conditions for the new cubical diagram
\msp{page 128}
(a) and c) moreover not having changed, and b) giving
\begin{itemize}
\item[b$'$)] $f'_1 : \cX/1 \longrightarrow \cX'/1$\ \ and\ \ $f'_2 : \cX/2 \longrightarrow \cX'/2$\ \ are in $W$.
\end{itemize}
\texttt{[rather} $f'_1 : 1\backslash\cX \longrightarrow 1\backslash\cX'$ \texttt{and} $f'_2 : 2\backslash\cX \longrightarrow 2\backslash\cX'$\texttt{]}). Then the conclusion of the corollary is equivalent to W(7) (dual formulation) in the case of the new cubical diagram, q.e.d.

\subsection*{3.2\quad Comparison of \ $\displaystyle\int
\begin{tikzcd}[column sep=small, row sep=tiny, ampersand replacement=\&, baseline=(\tikzcdmatrixname-2-1.base)]
 \& X_1 \\
X_0 \arrow[ur] \arrow[dr] \& \\
 \& X_2
\end{tikzcd}$ \ with an ordinary amalgamated sum.\\
$W$-homotopically cocartesian squares.}

The integrated category $\cX = \displaystyle\int
\begin{tikzcd}[column sep=small, row sep=tiny, ampersand replacement=\&, baseline=(\tikzcdmatrixname-2-1.base)]
 \& X_1 \\
X_0 \arrow[ur] \arrow[dr] \& \\
 \& X_2
\end{tikzcd}$ is a solution of the following universal problem, a 2-categorical variant of the one which defines the direct limit of diagrams. It is a matter, for $Y$ a variable object in Cat, of considering the systems of three functors $f_0$, $f_1$, $f_2$ and of the two functorial homomorphisms $\lambda_1$, $\lambda_2$:
\printedpage{87}
% head: A. Grothendieck : Dérivateurs, Ch. XII, Édition des Universités de Montpellier II et Paris VII
\[
\begin{tikzcd}[column sep=large, row sep=large, ampersand replacement=\&]
 \& X_0 \arrow[dl, "\alpha_1"'] \arrow[dr, "\alpha_2"] \arrow[dd, "f_0", ""{name=M, pos=0.62}] \& \\
X_1 \arrow[dr, "f_1"', ""{name=A, pos=0.45}] \& \& X_2 \arrow[dl, "f_2", ""{name=B, pos=0.45}] \\
 \& Y, \&
\arrow[from=A, to=M, bend left, shorten=3pt, "\lambda_1"]
\arrow[from=B, to=M, bend right, shorten=3pt, "\lambda_2"']
\end{tikzcd}
\tag*{$(*)$}
\]
\[
\begin{aligned}
\lambda_1 :\quad & f_1\alpha_1 \quad\longrightarrow\quad f_0\\
\lambda_2 :\quad & f_2\alpha_2 \quad\longrightarrow\quad f_0.
\end{aligned}
\]

\msp{page 129}
To give oneself an arrow of Cat
\[
\cX \longrightarrow Y
\]
``comes to the same thing'' as giving oneself a system $(f_0, f_1, f_2, \lambda_1, \lambda_2)$ as above, and this ``functorially in $Y$''. This is the special case of the general determination of the functors from the ``integrated'' fibered category of a contrafunctor
\[
I^{\circ} \longrightarrow \Cat, \qquad
\begin{aligned}
i &\quad\longmapsto\quad X_i\\
u &\quad\longmapsto\quad u^*
\end{aligned}
\]
to a category $Y$, in terms of a system of functors
\[
X_i \xrightarrow{f_i} Y,
\]
and of functorial homomorphisms
\[
f_i\circ u^* \xrightarrow{\lambda_u} f_j
\]
for $u : i \longrightarrow j$, whence $u^* : X_j \longrightarrow X_i$, with transitivity conditions for composite arrows, which we need not take into account in the present case, seeing that there are no composite arrows to look at in $I = \Phi$ (except for composites with an identity arrow, which we can dispense with looking at here \dots).

\msp{page 130}
If in particular we have a \textit{commutative} diagram in Cat

\end{document}
